\documentclass[11pt]{article}
\usepackage[a4paper,margin=1in]{geometry}
\usepackage[T1]{fontenc}
\usepackage[utf8]{inputenc}
\usepackage[english]{babel}
\usepackage{lmodern}
\usepackage{microtype}
\usepackage{amsmath,amssymb,mathtools}
\setlength{\parindent}{1.2em}
\setlength{\parskip}{0.35em}
\setlength{\emergencystretch}{4em}
\allowdisplaybreaks
\newcommand{\dd}{\,\mathrm d}
\newcommand{\cosec}{\operatorname{cosec}}
\newcommand{\art}[1]{\bigskip\begin{center}\large\bfseries #1.\end{center}\nopagebreak}
\newcommand{\tg}{\operatorname{tg}}
\newcommand{\arccosop}{\operatorname{arc\,cos}}
\begin{document}
\begin{center}
{\Large I.\\[0.4em] On the Elements of the Theory of the Eulerian Integrals.}\par
\medskip
{\normalsize [Inaugural dissertation, Göttingen 1852.]}
\end{center}

It is well known that the execution of the indirect operations in analysis usually meets with far greater difficulties than that of the direct operations; but precisely this seemingly unfortunate circumstance has always exercised the most favorable influence on the development of mathematics. Not only does overcoming these difficulties, when it is possible, always offer a peculiar attraction for the mathematician; rather, the very cases in which this could not be done with the concepts and aids formerly introduced have continually opened entirely new fields to the further development of mathematics. Thus, for example, the operations of subtraction, division, and extraction of roots lead to the concepts of negative, fractional, and imaginary numbers, each of which has enlarged the domain of mathematics in an extraordinary way. Very much the same is true in higher analysis with the underlying concept of function, which at first is based only on the operations taught in elementary mathematics (the corresponding functions could aptly be called elementary functions) and on their composition. The differential calculus, the most powerful lever for the development of the theory of functions, encounters no difficulties in its execution; that is, differentiating functions formed from the operations of elementary mathematics again leads to such functions. By contrast, the inverse kind of calculation, which in its entirety forms the integral calculus, has succeeded in doing the same only in comparatively few cases. In most cases it has not yet succeeded, or is perhaps entirely impossible, to represent the integrals of given functions with the aid of such functions. But precisely this circumstance has led to a considerable enlargement of the concept of function: new names and notations have been assigned to such non-representable integrals, and they have thereby been introduced into the circle of the earlier functions. In the development of the theory of such integral-functions, the cases in which they can be reduced to the functions hitherto alone in use are of the greatest importance, since their behavior is thereby made clearer and means are often also supplied for facilitating their computation. The chief purpose of the following treatise is to assemble these cases for the Eulerian integrals, with special attention to the method to be used.

\art{1}

It is first necessary briefly to recall the fundamental properties of the Eulerian integrals. The definitions of these functions lie in the equations
\begin{equation}
B(a,b)=\int_0^1 x^{a-1}(1-x)^{b-1}\dd x; \qquad
\Gamma(\mu)=\int_0^\infty x^{\mu-1}e^{-x}\dd x .
\tag{1}
\end{equation}
The integrals on the right are called Eulerian integrals of the first and of the second kind; the former are easily reduced to the latter. Namely, if in the double integral
\[
\int_0^\infty\!\int_0^\infty e^{-(x+y)}x^{a-1}y^{b-1}\dd y\dd x=\Gamma(a)\Gamma(b)
\]
one introduces two new variables $r$ and $w$ for $x$ and $y$ by setting $x+y=r$ and $x=rw$, whence $\dd y\dd x=r\dd r\dd w$, then one obtains
\[
\int_0^\infty e^{-r}r^{a+b-1}\dd r\int_0^1 w^{a-1}(1-w)^{b-1}\dd w=\Gamma(a)\Gamma(b),
\]
and hence, according to the definitions of $B$ and $\Gamma$,
\begin{equation}
B(a,b)=\frac{\Gamma(a)\Gamma(b)}{\Gamma(a+b)},
\tag{2}
\end{equation}
from which it follows at the same time that $B(b,a)=B(a,b)$ is a symmetric function of $a$ and $b$. This can also be shown directly by putting $1-x$ for $x$ in the integral $B$.

If in the integral $\Gamma$ one puts $rx$ for $x$ and takes $r$ as a positive constant, then one obtains
\begin{equation}
\int_0^\infty x^{\mu-1}e^{-rx}\dd x=\frac{\Gamma(\mu)}{r^\mu},
\tag{3}
\end{equation}
and from this, by $n$-fold partial differentiation with respect to $r$, the important relation
\begin{equation}
\Gamma(\mu+n)=(\mu+n-1)(\mu+n-2)\cdots(\mu+1)\mu\,\Gamma(\mu).
\tag{4}
\end{equation}
It is clear that it is necessary to compute $\Gamma(\mu)$ only for each $\mu$ within an interval of length one in order, with the aid of this equation, to find $\Gamma(\mu)$ for every other $\mu$.

\art{2}

Important consequences for the theory of Eulerian integrals may be drawn from the formulas just developed, especially with respect to the cases in which they can be represented without the aid of new functions. Since those of the first kind can be reduced to those of the second, we begin with the investigation of the latter. It is known that definite integrals can be found whenever the indefinite integrals can be represented generally (cf. Art. 6); the integral calculus teaches, however, that this is possible for the integral
\[
\int x^{\mu-1}e^{-x}\dd x
\]
only when $\mu$ is a positive integer $n$. We may therefore conclude in advance that then $\Gamma(n)$ also can be given. But in equation (4) we have obtained a reduction formula which teaches us how one gamma-function can be represented by another when their arguments differ by an integer. Taking, therefore, most simply $\mu=1$, we obtain from the indefinite integration
\begin{equation}
\Gamma(1)=\int_0^\infty e^{-x}\dd x=1 \quad\text{and consequently}\quad
\Gamma(n)=(n-1)(n-2)\cdots2\cdot1.
\tag{5}
\end{equation}
This, however, is also the only case in which it is evident a priori that $\Gamma(\mu)$ can be represented.

We can take a similar path with the Eulerian integrals of the first kind, by first examining the representability of the indefinite integral
\[
\int x^{a-1}(1-x)^{b-1}\dd x.
\]
This belongs, as is known, to the class of integrals of so-called binomial differentials, which stand under the general form
\[
\int x^m(a+bx^n)^p\dd x.
\]
It is known that the binomial differential can always be made rational, and hence integrable, when either
\[
\frac{m+1}{n}\quad\text{or}\quad \frac{m+1}{n}+p
\]
is an integer, and when $m,n,p$ are rational. Thus, in order that the above indefinite integral be representable, either $a$ [or also $b$, since $B(a,b)=B(b,a)$] or $a+b$ must be an integer. But the first case also follows immediately from formulas (2) and (4). For if $a$ is an integer $m$, these formulas give
\[
B(m,b)=\frac{\Gamma(m)\Gamma(b)}{\Gamma(m+b)}=
\frac{\Gamma(m)}{(m-1+b)(m-2+b)\cdots(1+b)b},
\]
and hence, with the aid of formula (5),
\begin{equation}
B(m,b)=\frac{(m-1)(m-2)\cdots2\cdot1}{(m-1+b)(m-2+b)\cdots(1+b)b}.
\tag{6}
\end{equation}
Likewise, if $n$ denotes a positive integer, one obtains
\[
B(a,n)=\frac{(n-1)(n-2)\cdots2\cdot1}{(n-1+a)(n-2+a)\cdots(1+a)a},
\]
\[
B(m,n)=\frac{(m-1)\cdots2\cdot1\,(n-1)\cdots2\cdot1}{(m+n-1)(m+n-2)\cdots2\cdot1}.
\]
By contrast, the second case, in which $a+b$ is an integer without $a$ and $b$ both being integers, supplies, independently of the Eulerian integrals of the second kind, a new class of integrals of which one can assert a priori that they can be represented. Nevertheless they can all be reduced to a single one as follows. If $a+b$ is a positive integer (positive, because it is to be assumed known that the Eulerian integrals for negative arguments always become infinitely large), one can always put $a=m+r$, $b=n-r$, where $m$ and $n$ are positive integers and $r$ is a positive proper fraction. Using the reduction formula (4), one then easily finds
\[
B(m+r,n-r)=\frac{\Gamma(m+r)\Gamma(n-r)}{\Gamma(m+n)}
\]
\[
=\frac{(m-1+r)(m-2+r)\cdots r\Gamma(r)\,(n-1-r)(n-2-r)\cdots(1-r)\Gamma(1-r)}{(m+n-1)(m+n-2)\cdots3\cdot2\cdot1},
\]
or, since from equations (2) and (5) $\Gamma(r)\Gamma(1-r)=B(r,1-r)$,
\begin{equation}
B(m+r,n-r)=\frac{(m-1+r)\cdots r\,(n-1-r)\cdots(1-r)}{(m+n-1)(m+n-2)\cdots2\cdot1}\,B(r,1-r).
\tag{7}
\end{equation}
The integral $B(r,1-r)$, to which the integrals $B(m+r,n-r)$ are thereby reduced, is one of the most interesting in the integral calculus and is especially of the greatest importance for the theory of Eulerian integrals of both kinds, as is already clear from the last equation. As simple as the form is in which its value can be represented, the methods by which that value is found are mutually diverse and complicated. The following articles will therefore be concerned with this integral.

\art{3}

For the sake of uniformity of notation I shall write $b$ instead of $r$, and denote $B(b,1-b)=\Gamma(b)\Gamma(1-b)$ briefly by $B$, so that $B$ is regarded as a function of the single variable $b$. Thus the following equation holds between $b$ and $B$:
\begin{equation}
\left\{
\begin{aligned}
B&=\int_0^1 \left(\frac{x}{1-x}\right)^b\frac{\dd x}{x}
  =\int_0^1 \left(\frac{x}{1-x}\right)^{1-b}\frac{\dd x}{x} \\
&=\int_0^\infty \frac{x^{b-1}\dd x}{x+1}
  =\int_0^\infty \frac{x^{-b}\dd x}{x+1}.
\end{aligned}
\right.
\tag{8}
\end{equation}
The last two forms are obtained when one writes $x$ for $x/(1-x)$ in the first two. Although in the preceding article $b$ is already restricted to the interval between $0$ and $1$, because otherwise the integral $B$, as an Eulerian integral, would have a negative argument and hence an infinitely large value, the necessity of this restriction shall now be proved in all rigor, independently of the theory of Eulerian integrals. For greater generality I do not wish to attach this investigation directly to the integral $B$, but to the more general integral
\[
\int_0^\infty \frac{1+x+xx+\cdots+x^{m-1}}{1+x+xx+\cdots+x^{n-1}}x^{b-1}\dd x
=\int_0^\infty \frac{1-x^m}{1-x^n}x^{b-1}\dd x,
\]
which, for $m=1$, $n=2$, passes into the integral $B$. The form on the right does not even require that $m$ and $n$ be integers. I shall therefore pass immediately to the consideration of the integral
\begin{equation}
\varphi(b)=\int_0^\infty \frac{1-x^\mu}{1-x^\nu}x^{b-1}\dd x,
\tag{9}
\end{equation}
in which $b,\mu,\nu$ are to denote arbitrary real constants.

First, one readily sees that $\mu$ and $\nu$ may always be assumed positive without impairing the generality of the integral, since instead of $1-x^\mu$ and $1-x^\nu$ one may also write $x^\mu(x^{-\mu}-1)$ and $x^\nu(x^{-\nu}-1)$. If one then decomposes $\varphi(b)$ into two integrals of the same function, with limits respectively $0,1$ and $1,\infty$, and in the second integral puts $1/x$ for $x$, then this integral also obtains the limits $0,1$, and both can be combined as
\begin{equation}
\varphi(b)=\int_0^1 \frac{1-x^\mu}{1-x^\nu}\left(x^b+x^{\nu-\mu-b}\right)\frac{\dd x}{x}.
\tag{10}
\end{equation}
Since, by assumption, $\mu$ and $\nu$ are positive, within the limits of integration for $x$, that is, when $x$ is a positive proper fraction, the quotient
\[
\frac{1-x^\mu}{1-x^\nu}
\]
is throughout a positive finite number; for also at $x=1$ this quotient obtains a finite value, namely $\mu/\nu$. If its greatest and least values are denoted by $M$ and $N$, these are likewise finite positive numbers, and throughout the interval of integration
\[
M>\frac{1-x^\mu}{1-x^\nu}>N.
\]
Consequently, by a known theorem from the theory of definite integrals,
\[
M\int_0^1\left(x^b+x^{\nu-\mu-b}\right)\frac{\dd x}{x}
>\varphi(b)>
N\int_0^1\left(x^b+x^{\nu-\mu-b}\right)\frac{\dd x}{x}.
\]
Now the integral which occurs on both sides of $\varphi(b)$ with the factors $M$ and $N$ has the values
\[
\frac{\nu-\mu}{b(\nu-\mu-b)},\quad \pm\infty, \quad\text{or}\quad -\infty,
\]
according as both $b$ and $\nu-\mu-b$ are positive, or one of them is zero, or one of them is negative. It follows that only in the first case does $\varphi(b)$ have a finite, and indeed positive, value; in every other case it has an infinitely large value. It is therefore necessary that both $b$ and $\nu-\mu-b$ be positive numbers, which can be expressed by the condition $\nu-\mu>b>0$. From this it follows at the same time that $\nu>\mu$. If $b=0$ or $b=\nu-\mu$, then $\varphi(b)$ becomes not only infinitely large but also discontinuous, jumping from $+\infty$ to $-\infty$.

A further interesting consequence can be drawn from equation (10). If one puts $\nu-\mu-b$ in place of $b$, one immediately obtains
\begin{equation}
\varphi(\nu-\mu-b)=\varphi(b).
\tag{11}
\end{equation}
If in this one sets $b=(\nu-\mu)/2+b'$, differentiates with respect to $b'$, and then sets $b'=0$, it follows that $\varphi'((\nu-\mu)/2)=0$, where $\varphi'(b)=\dd\varphi(b)/\dd b$. To decide whether the value $b=(\nu-\mu)/2$ corresponds to a maximum or a minimum of $\varphi(b)$, one must form the second differential quotient $\varphi''(b)$; since this is represented by the integral
\[
\int_0^\infty \frac{1-x^\mu}{1-x^\nu}x^{b-1}\dd x\,(\ell x)^2,
\]
where $\ell x$ denotes the natural logarithm of $x$, and is consequently positive throughout the whole interval of $b$, $\varphi((\nu-\mu)/2)$ is a minimum of $\varphi(b)$, and indeed the only one. This minimum is therefore
\begin{equation}
\varphi\left(\frac{\nu}{2}-\frac{\mu}{2}\right)=
\int_0^\infty \frac{x^{\mu/2}-x^{-\mu/2}}{x^{\nu/2}-x^{-\nu/2}}\frac{\dd x}{x}
=2\int_0^\infty \frac{x^\mu-x^{-\mu}}{x^\nu-x^{-\nu}}\frac{\dd x}{x}.
\tag{12}
\end{equation}
Applying the preceding to our case, in which $\mu=1$, $\nu=2$, one obtains that the integral $B$ has a finite, positive value only when $b$ is a positive proper fraction. Furthermore $B=\varphi(b)=\varphi(1-b)$, and for $b=\tfrac12$ it reaches the minimum
\begin{equation}
\varphi\left(\frac12\right)=2\int_0^\infty \frac{x-1/x}{xx-1/(xx)}\frac{\dd x}{x}
=2\int_0^\infty \frac{\dd x}{xx+1}=\pi.
\tag{13}
\end{equation}

\art{4}

We shall now give the principal proofs that have so far been established for the value of the integral $B$. It will suffice, however, briefly to indicate their course and to enter into detail only where a stricter foundation seems necessary.

Since Art. 2 has already shown that the value of the integral $B(m+r,n-r)$, of which $B$ is only a special case, must result from indefinite integration when $r$ is a proper rational fraction, it is most natural to begin with this method. Thus take $b=m/n$, where $m$ and $n$ are positive integers and $m<n$. The first matter is then to carry out the indefinite integration in the integral
\begin{equation}
B=\int_0^\infty \frac{x^{b-1}\dd x}{x+1}
=\int_0^\infty \frac{x^{m/n-1}\dd x}{x+1}
=n\int_0^\infty \frac{x^{m-1}\dd x}{x^n+1},
\tag{14}
\end{equation}
which, as is known, requires decomposing the function under the integral sign into partial fractions. Put, for abbreviation, $\vartheta=e^{\pi i/n}$, where $i=\sqrt{-1}$. Then
\[
x^n+1=(x-\vartheta)(x-\vartheta^3)\cdots(x-\vartheta^{2k-1})\cdots(x-\vartheta^{2n-1}).
\]
If one then carries out the decomposition into partial fractions with linear denominators and integrates each individual term, one finds
\begin{equation}
n\int \frac{x^{m-1}\dd x}{x^n+1}
=-\sum_{k=1}^n \vartheta^{m(2k-1)}\ell(\vartheta^{2k-1}-x).
\tag{15}
\end{equation}
The sum on the right can also be written as
\[
-\sum_{k=1}^n \vartheta^{m(2k-1)}\ell\left(\frac{\vartheta^{2k-1}}{x}-1\right)
-\ell x\sum_{k=1}^n\vartheta^{m(2k-1)},
\]
where
\[
\sum\vartheta^{m(2k-1)}=\vartheta^m\sum(\vartheta^{2m})^{k-1}
=\vartheta^m\frac{\vartheta^{2mn}-1}{\vartheta^{2m}-1}=0;
\]
hence also
\begin{equation}
n\int \frac{x^{m-1}\dd x}{x^n+1}
=-\sum_{k=1}^n\vartheta^{m(2k-1)}\ell\left(\frac{\vartheta^{2k-1}}{x}-1\right).
\tag{16}
\end{equation}
It should be carefully noted here that the sums in (15) and (16) are completely equal, not merely different by a constant. If, therefore, one puts $x=\infty$ in (16) and $x=0$ in (15), the difference gives the definite integral $B$, namely
\begin{equation}
\left\{
\begin{aligned}
B&=\sum_{k=1}^n\vartheta^{m(2k-1)}\ell(\vartheta^{2k-1})\\
&=\frac{\pi i}{n}\sum_{k=1}^n(2k-1)\vartheta^{m(2k-1)}=\frac{\pi}{\sin b\pi}.
\end{aligned}
\right.
\tag{17}
\end{equation}
This gives the value of $B$ for rational $b$. By transforming (15) into (16) I believe I have shown in the shortest way that the sum in (15) vanishes for $x=\infty$; and with respect to rigor, this path seems at least not inferior to those followed in most textbooks of the integral calculus.

\art{5}

A second proof is the following, which, so far as I know, was given by Schlömilch. From the equation
\[
B=\int_0^1\frac{x^{b-1}+x^{-b}}{x+1}\dd x,
\]
which follows from formula (10) in Art. 3 by putting $\mu=1$, $\nu=2$, one obtains, by expanding $1/(x+1)$ in powers of $x$ with regard to the remainder and by carrying out the integrations for $B$, the $n$-term series
\begin{equation}
\frac1b+\frac{2b}{1-b^2}-\frac{2b}{4-b^2}+\frac{2b}{9-b^2}-\cdots+(-1)^n\frac{2b}{(n-1)^2-b^2}
\tag{18}
\end{equation}
together with the remainder
\[
\frac{(-1)^{n-1}}{n-b}+(-1)^n\int_0^1\frac{(x^{b-1}+x^{-b})x^n}{x+1}\dd x.
\]
By arguments similar to those made in Art. 3 concerning the finiteness of $\varphi(b)$, it can be shown that this remainder tends to zero as $n$ increases without bound, so that $B$ is to be set equal to the infinitely continued series (18). But if in the known formula
\[
\cosec u=\frac1u+\frac{2u}{\pi^2-u^2}+\frac{2u}{4\pi^2-u^2}+\text{etc.}
\]
one takes $u=b\pi$, the same value for $B$ is found immediately as in the preceding article.

\art{6}

A third proof, due to Cauchy, rests on a theorem concerning the so-called reversal of the order of integration in double integrals. Since the considerations belonging here are very delicate, and since obscurities are still often found concerning them, I ask leave to go somewhat further back in order to obtain a secure foundation for this investigation. For this it is necessary to return to the foundations of the theory of definite integrals.

The definite integral is usually defined as the difference of two values of the indefinite integral corresponding to two special values of the variable of integration; the latter are called the limits of the integral. Simple as this definition seems, it cannot satisfy stricter demands, which must always be made when one is fixing a basis for a whole theory. The chief reason for rejecting this definition lies especially in the circumstance that it is not based directly on the function actually given, but presupposes as an intermediary another function, namely the indefinite integral; and this is a defect in several respects. First, the existence of the definite integral is not evident until that of the indefinite integral has been proved. Even if this were generally possible, one must ask, on the other hand, whether according to this definition the definite integral is really to be called definite, that is, whether it depends only on the given function and the limits. It has already been shown several times that this is by no means the case, and cases are known in which this definition leads to ambiguities that cannot be removed by this route alone. I hope now to show that the definite integral understood according to this definition is in every case just as indeterminate as the so-called indefinite integral.

While the above definition already gives occasion for doubt when it is not possible, with the aid of the known methods, to represent the indefinite integral, this happens to a much greater degree when there are several, indeed infinitely many, functions whose differential is the given function. It can be rigorously proved that these functions can differ from one another only by so-called constants, and it is precisely on this that the above definition rests, since it tacitly presupposes that the constant difference of such functions remains truly the same for all values of the variable. But precisely this is by no means necessary. One may instead conceive that this constant has different values in different finite intervals of the variable $x$, and yet in each of them the differential quotient of $f(x)+C$ will be the same function $f'(x)$, provided that $C$ does not vary continuously with $x$, for then $C$ would no longer be a constant. This is especially clear geometrically if $f(x)$ is regarded as the ordinate of a curve: arbitrary pieces of this curve can be displaced parallel to the axis of ordinates without thereby changing $f'(x)$. In other words, the first differential quotient of a function gives not the slightest information about the continuity of the function. Such discontinuities can also be represented analytically, notably with the aid of Fourier integrals, but still more simply with wholly elementary means.

Because of the ambiguity that attaches to every square root with respect to its sign, it is absolutely necessary always to denote only one root by a definite sign. Thus it has also been agreed that $\sqrt{x}$ is always to mean the positive square root of $x$. The necessity of this is especially evident when in place of a radical expression another expression is put, for example the binomial series, which in any event expresses only one root. Given this, one can easily form an expression that does not change continuously with $x$, and is therefore a constant, but nevertheless takes different values in different intervals. Such an expression is, for example,
\[
C\frac{x-c}{\sqrt{(x-c)^2}},
\]
which is equal to $+C$ or $-C$ according as $x$ is taken greater or less than $c$. By differentiation one of course finds
\[
\frac{\dd}{\dd x}\left(C\frac{x-c}{\sqrt{(x-c)^2}}\right)
=C\frac{\sqrt{(x-c)^2}-(x-c)\frac{x-c}{\sqrt{(x-c)^2}}}{(x-c)^2}=0
\]
and consequently also
\[
\frac{\dd}{\dd x}\left(f(x)+C\frac{x-c}{\sqrt{(x-c)^2}}\right)=\frac{\dd f(x)}{\dd x}=f'(x).
\]
If, therefore, one wanted to apply the above definition of the definite integral, one would have to obtain from the indefinite integral
\[
\int f'(x)\dd x=f(x)+C\frac{x-c}{\sqrt{(x-c)^2}}
\]
the definite integral
\[
\int_a^b f(x)\dd x=f(b)-f(a)+C\frac{b-c}{\sqrt{(b-c)^2}}-C\frac{a-c}{\sqrt{(a-c)^2}}.
\]
If $a<c<b$, the sum of the last two terms is equal to $2C$, so that the definite integral still contains a wholly arbitrary constant.

This sufficiently exhibits the inadequacy of this definition of the indefinite integral, and we turn now to another definition that starts directly from the given quantities. According to it, the definite integral is to be understood as the sum of all the infinitely small values of the given differential when the variable $x$ is assigned, continuously, all values lying between the given limits of the integral. This definition permits no ambiguities other than those already lying in the nature of the given function. It can further be shown that it is always identical with the former definition, provided only that the indefinite integral is chosen so that it contains no discontinuity within the interval of integration. For then the sum which, according to the second definition, forms the integral is precisely the sum of all the infinitely small increments that the indefinite integral $f(x)$ receives as $x$ passes continuously through the interval from $a$ to $b$. If, however, $f(x)$ is discontinuous at some point $c$ between $a$ and $b$, one may imagine substituting a continuous function $f_1(x)$ whose differential is likewise $f'(x)\dd x$. Then the definite integral is $f_1(b)-f_1(a)$, and this difference differs from $f(b)-f(a)$ only by the amount of the jump that $f(x)$ makes at the point $x=c$; hence
\[
\int_a^b f'(x)\dd x=f(b)-f(a)+\lim\bigl(f(c-\delta)-f(c+\varepsilon)\bigr),
\]
where $\delta$ and $\varepsilon$ denote infinitely small quantities having the same sign as $(b-a)$. Instead of this equation one can also write
\[
\int_a^b f'(x)\dd x=\lim\left(\int_a^{c-\delta}f'(x)\dd x+\int_{c+\varepsilon}^b f'(x)\dd x\right).
\]

\art{7}

The preceding considerations are particularly important for such definite integrals as contain still another variable, that is, for integrals of the form
\[
\int_a^b f(x,\xi)\dd x,
\]
where $\xi$ denotes a variable independent of $x$. It may happen that the indefinite integral, although in general it is continuous for all values of $x$, loses this property when a definite value is assigned to the variable $\xi$. This is to be considered especially when the definite integral, regarded as a function of $\xi$, is subjected to a second integration with respect to $\xi$, and indeed between limits within which also lies the special value of $\xi$ that makes the indefinite integral taken with respect to $x$ discontinuous.

Let $\alpha,\beta$ be these limits, and $c$ and $\gamma$ the values of $x$ and $\xi$ for which the integral with respect to $x$ becomes discontinuous. Then, according to the preceding article, one must set
\begin{equation}
\int_\alpha^\beta \dd\xi\int_a^b f(x,\xi)\dd x=
\lim\int_\alpha^\beta\dd\xi\int_a^{c-\varepsilon} f(x,\xi)\dd x+
\int_\alpha^\beta\dd\xi\int_{c+\varepsilon}^b f(x,\xi)\dd x.
\tag{19}
\end{equation}
For the integration with respect to $x$ must first be carried out in such a way that it remains valid for all values of $\xi$ that enter into consideration in the second integration. I have cited this formula in order to counter the erroneous interpretation of a theorem that concerns the reversal of the order of integration in double integrals. In such a double integral one may indeed interchange the order, and hence put
\begin{equation}
\int_\alpha^\beta \dd\xi\int_a^b f(x,\xi)\dd x=
\int_a^b \dd x\int_\alpha^\beta f(x,\xi)\dd\xi,
\tag{20}
\end{equation}
if $f(x,\xi)$ remains finite and continuous for all values of $x$ and $\xi$ within the integration. But if $f(x,\xi)$ becomes discontinuous for $x=c$, $\xi=\gamma$, one can still put
\[
\int_\alpha^\beta\dd\xi\int_a^{c-\varepsilon}f(x,\xi)\dd x
+\int_\alpha^\beta\dd\xi\int_{c+\varepsilon}^b f(x,\xi)\dd x
=\int_a^{c-\varepsilon}\dd x\int_\alpha^\beta f(x,\xi)\dd\xi
+\int_{c+\varepsilon}^b\dd x\int_\alpha^\beta f(x,\xi)\dd\xi.
\]
If the indefinite integrals with respect to $x$ and $\xi$ are denoted respectively by $F(x,\xi)$ and $\varphi(x,\xi)$ and are assumed continuous between the limits of the individual integrals, the last equation becomes
\[
\int_\alpha^\beta [F(b,\xi)-F(a,\xi)]\dd\xi-
\int_\alpha^\beta [F(c+\varepsilon,\xi)-F(c-\varepsilon,\xi)]\dd\xi
\]
\[
=\int_a^{c-\varepsilon}[\varphi(x,\beta)-\varphi(x,\alpha)]\dd x
+\int_{c+\varepsilon}^b[\varphi(x,\beta)-\varphi(x,\alpha)]\dd x.
\]
Letting $\varepsilon$ become zero, one obtains
\begin{equation}
\left\{
\begin{aligned}
&\int_\alpha^\beta [F(b,\xi)-F(a,\xi)]\dd\xi
-\lim\int_\alpha^\beta [F(c+\varepsilon,\xi)-F(c-\varepsilon,\xi)]\dd\xi \\
&\hspace{7em}=\int_a^b[\varphi(x,\beta)-\varphi(x,\alpha)]\dd x .
\end{aligned}
\right.
\tag{21}
\end{equation}
This formula is usually understood as if the second term on the left side gave the difference between the two double integrals in (20). But what was said at the beginning of this article shows that this is not correct, since only both terms of the left side taken together represent the double integral on the left side of (20). This does not impair the correctness of equation (21), and it is precisely on this equation that Cauchy's proof rests.

Namely, take
\[
f(x,\xi)=i f'(x+\xi i), \qquad i=\sqrt{-1}, \qquad f'(z)=\frac{\dd f(z)}{\dd z}.
\]
Then
\[
F(x,\xi)=i f(x+\xi i),\qquad \varphi(x,\xi)=f(x+\xi i),
\]
and equation (21) passes into the following:
\begin{equation}
\left\{
\begin{aligned}
&i\int_\alpha^\beta [f(b+\xi i)-f(a+\xi i)]\dd\xi \\
&\quad -i\lim\int_\alpha^\beta [f(c+\varepsilon+\xi i)-f(c-\varepsilon+\xi i)]\dd\xi \\
&\hspace{8em}=\int_a^b [f(x+\beta i)-f(x+\alpha i)]\dd x .
\end{aligned}
\right.
\tag{22}
\end{equation}
If in the second term on the left one introduces a new variable $\eta$ by the equation $\xi=\gamma+\varepsilon\eta$, where by assumption $\alpha<\gamma<\beta$, and sets
\begin{equation}
f(z)=\frac{F(z)}{z-c-\gamma i},
\tag{23}
\end{equation}
one easily finds
\begin{equation}
i\lim\int_\alpha^\beta [f(c+\varepsilon+\xi i)-f(c-\varepsilon+\xi i)]\dd\xi=2\pi i F(c+\gamma i).
\tag{24}
\end{equation}

\art{8}

From the formulas just developed Cauchy derived the value of the integral $B$, but in a way that in individual points appears much in need of a stricter justification. It consists in the following. If the function $f(z)$ is so constituted that for every value of $\xi$, $f(\pm\infty+\xi i)=0$, and for every value of $x$, $f(x+\infty i)=0$, then from equations (22), (23), and (24), by putting $\alpha=0$, $\beta=\infty$, $a=-\infty$, $b=\infty$, it follows that
\begin{equation}
\int_{-\infty}^{+\infty} f(x)\dd x=2\pi i F(c+\gamma i),
\tag{25}
\end{equation}
where now $\gamma$ must necessarily be positive by the condition $\alpha<\gamma<\beta$. If one now sets
\[
f(z)=\frac{(-zi)^{\mu-1}}{zz+1},
\]
where $\mu$ is a number lying between $0$ and $2$ (without motivation), then the conditions $f(\pm\infty+\xi i)=0$ and $f(x+\infty i)=0$ are satisfied. The values of $x$ and $\xi$ that make $f(x+\xi i)$ infinite are $c=0$, $\gamma=1$; hence
\[
F(z)=(z-i)\frac{(-zi)^{\mu-1}}{zz+1},\qquad F(c+\gamma i)=F(i)=\frac1{2i},
\]
and consequently
\begin{equation}
\int_{-\infty}^{+\infty}\frac{(-xi)^{\mu-1}}{xx+1}\dd x=\pi.
\tag{26}
\end{equation}
If this integral is decomposed into two others with limits $0,\infty$ and $-\infty,0$, and in the second one puts $-x$ for $x$, one finds
\[
\int_0^\infty\frac{x^{\mu-1}\dd x}{xx+1}
=\frac{\pi}{(+i)^{\mu-1}+(-i)^{\mu-1}}
=\frac{\pi}{2\cos(\mu-1)\frac{\pi}{2}}
=\frac{\pi}{2\sin\mu\frac{\pi}{2}}.
\]
It remains only to put $x$ for $xx$ and $\mu=2b$ (where $b$ lies between $0$ and $1$ when $0<\mu<2$) in order to recover equation (17).

In this derivation it is especially the derivation of equation (25) that does not seem to me entirely rigorous. For although the conditions $f(\pm\infty+\xi i)=0$, $f(x+\infty i)=0$ are satisfied for every $\mu$ lying between $0$ and $2$ (indeed for this it is only necessary that $\mu<3$, so that $\mu$ may also be negative), it is nevertheless known that the vanishing of the function under the integral sign does not always imply the vanishing of the integral, especially when one of the limits is infinitely large. I shall therefore try, through the following presentation, to remove these doubts and at the same time to prove that $\mu$ must lie between the limits $0$ and $2$.

Put in equation (22) $\alpha=0$, $b=\beta=-a=k$. Taking equation (24) into account, it becomes
\[
i\int_0^k [f(k+\xi i)-f(-k+\xi i)]\dd\xi-2\pi i F(c+\gamma i)
=\int_{-k}^{+k} f(x+ki)\dd x-\int_{-k}^{+k} f(x)\dd x.
\]
If in the first integral one puts $\xi=k\eta$, and in the second $x=ky$, one obtains
\[
i\int_0^1 [f(k(1+\eta i))-f(-k(1-\eta i))]k\dd\eta-2\pi iF(c+\gamma i)
=\int_{-1}^{+1}f(k(y+i))k\dd y-\int_{-k}^{+k}f(x)\dd x.
\]
If now, as $k$ increases without bound, the functions under the integral signs vanish, and indeed for every value of the variables, then the integrals themselves become zero. In our case
\[
kf(k(1+\eta i))=(-i)^{\mu-1}\frac{k^\mu(1+\eta i)^{\mu-1}}{kk(1+\eta i)^2+1},
\]
and the other functions are similar. In order that these expressions vanish as $k$ increases without bound, it is necessary that $\mu<2$; this does not exclude, however, that $\mu$ may also be negative. Under this assumption, in any case, one obtains equation (25). But from the course of the proof in the preceding article it is clear that if there are several pairs of values, such as $c$ and $\gamma$, for which $f(x,\xi)$ becomes infinite, then in equation (25) the sum of the corresponding expressions must be taken (with only the remark that, if $\gamma=\alpha$, in equation (24) $\pi iF(c+\gamma i)$ must be put instead of $2\pi iF(c+\gamma i)$). In our case, however, $f(x,\xi)=i f'(x+\xi i)$, and according to the above specialization of $f(z)$,
\[
f'(z)=(-i)^{\mu-1}\frac{(\mu-3)z^\mu+(\mu-1)z^{\mu-2}}{(zz+1)^2}.
\]
One must seek the complex values of $z$ that make this function infinite. The first are obtained from the equation $zz+1=0$, giving the two systems $(c=0,\gamma=1)$ and $(c=0,\gamma=-1)$. The first has already been treated above; the second must be excluded because this $\gamma$ does not satisfy the condition $\alpha<\gamma<\beta$. But if, as just shown, $\mu<2$ must hold, then $(c=0,\gamma=0)$ is also such a system, and we would therefore have to add the correction $\pi iF(0)$, where $F(z)=(z-c-\gamma i)f(z)$, hence in this case
\[
F(z)=(-i)^{\mu-1}\frac{z^\mu}{zz+1}.
\]
For $z=0$, $F(z)$ is either infinitely large or zero, according as $\mu$ is negative or positive. Thus, if the value of the integral in (25) is to be finite, we must assume the latter; then no correction is to be added to (26), since $F(0)=0$. By this consideration the validity of equation (26), from which the value of the integral $B$ follows immediately, is therefore restricted to the condition $0<\mu<2$, or $0<b<1$.

\art{9}

A proof essentially different from those so far cited is finally contained in the treatise “Disquisitiones generales circa seriem infinitam etc. Auctore C. F. Gauss.” The whole arrangement of that work, however, makes it impossible to present the proof here, since the foundations on which it rests yield at once a complete theory of the Eulerian integrals, and for that very reason are too substantial to be developed here merely for this single purpose.

To these proofs I now wish to add one which, so far as I know, is new, which requires very little preparation, and which always remains within the domain of integral calculus, even though the method itself is not exactly new: it amounts to reducing the problem to the integration of a differential equation of the second order. In the equation
\[
B B=\int_0^\infty\frac{x^{b-1}\dd x}{x+1}
     \int_0^\infty\frac{y^{b-1}\dd y}{y+1}
=\int_0^\infty\frac{\dd x}{x+1}
     \int_0^\infty\frac{(xy)^{b-1}}{y+1}\dd y
\]
put $y=\frac zx$, $\dd y=\frac{\dd z}{x}$; then reverse the order of integration and carry out the integration with respect to $x$. One then easily finds
\begin{equation}
B B=\int_0^\infty\frac{z^{b-1}\ell z}{z-1}\dd z
=\frac{\dd}{\dd b}\int_0^\infty\frac{z^{b-1}\dd z}{z-1}.
\tag{27}
\end{equation}
By indefinite integration with respect to $b$ one therefore obtains
\[
\int B B\,\dd b=\int_0^\infty\frac{z^{b-1}\dd z}{z-1}.
\]
The integral on the right differs from the integral $B$ only in the form of the denominator, and it is easy to foresee that one recovers the integral $B$ if the expression just obtained is subjected to the same treatment. Indeed, one finds
\[
B\int B B\,\dd b
=\int_0^\infty\frac{\dd z}{z-1}
  \int_0^\infty\frac{(zy)^{b-1}}{y+1}\dd y.
\]
If one now puts $y=\frac xz$, $\dd y=\frac{\dd x}{z}$, then, according to Articles 6 and 7, decomposes the integral taken with respect to $z$ into two integrals whose limits are $0,1-\delta$ and $1+\varepsilon,\infty$, reverses the order of integration, and carries out the integration with respect to $z$, one obtains
\[
B\int B B\,\dd b
=\int_0^\infty\frac{x^{b-1}\ell x}{x+1}\dd x
+\lim\int_0^\infty\frac{x^{b-1}\dd x}{x+1}\,\ell\!\left(\frac{\delta}{\varepsilon}\right).
\]
Here $\lim \ell(\delta/\varepsilon)$ is admittedly undetermined, but in any case it is independent of $x$ and $b$; let it be denoted by $k$. Then the last equation gives
\[
B\int B B\,\dd b=\frac{\dd B}{\dd b}+kB,
\]
or, if one observes that the integral on the left already contains an arbitrary constant,
\begin{equation}
B\int B B\,\dd b=\frac{\dd B}{\dd b}.
\tag{28}
\end{equation}
From this, by division by $B$ and differentiation with respect to $b$, one obtains the second-order differential equation
\[
B B=\frac1B\frac{\dd\dd B}{\dd b^2}
-\frac1{B B}\left(\frac{\dd B}{\dd b}\right)^2.
\]
It has the distinctive property that the independent variable $b$ does not occur; it can therefore, as is well known, be reduced to a differential equation of first order if the first differential quotient is introduced as a new variable. Denote this quotient by $B'$. Then
\[
\frac{\dd B}{\dd b}=B',\qquad
\frac{\dd\dd B}{\dd b^2}=\frac{\dd B'}{\dd b}
=\frac{B'\dd B'}{\dd B}.
\]
If these transformations are introduced into the above differential equation, it becomes
\[
B B=\frac{B'\dd B'}{B\,\dd B}-\frac{B'B'}{B B},
\]
or
\[
\dd(BB)=
\frac{BB\,\dd(B'B')-B'B'\dd(BB)}{(BB)^2}
=\dd\frac{B'B'}{BB}.
\]
Its integral is
\[
BB=cc+\frac{B'B'}{BB}
=cc+\frac1{BB}\left(\frac{\dd B}{\dd b}\right)^2.
\]
To determine the constant, it is now necessary to know two properties of the integral. We take these from Article 3, where it was shown that for $b=\frac12$ the integral $B$ attains a minimum equal to $\pi$; since at the same time $\frac{\dd B}{\dd b}=0$, it follows immediately that $cc=\pi\pi$.

One then obtains further
\[
\pm\,\dd b=
\frac{\dd B}{B\sqrt{BB-\pi\pi}}
=\frac1\pi
  \frac{-\dd\frac{\pi}{B}}
       {\sqrt{1-\frac{\pi\pi}{BB}}},
\]
and therefore, by integration,
\[
\pm(b+c)=\frac1\pi\,\arccosop\frac{\pi}{B},
\qquad
B=\frac{\pi}{\cos(b+c)\pi}.
\]
Since $B$ has the value $\pi$ for $b=\frac12$, one must have
\[
\cos\left(\frac12+c\right)\pi=-\sin c\pi=1,
\qquad \cos c\pi=0.
\]
Consequently
\[
\cos(b+c)\pi=\cos b\pi\cos c\pi-\sin b\pi\sin c\pi=\sin b\pi,
\]
whereby one again finds $B=\frac{\pi}{\sin b\pi}$.

\art{10}

Because the proof given in the preceding article still does not fully satisfy the requirements of the greatest rigor, especially insofar as the integral
\[
\int_0^\infty\frac{x^{b-1}\dd x}{x-1}
\]
plays a role in it, it is not inappropriate to add here a modification in which indefinite integration with respect to $b$ is entirely avoided. Since equation (28) can also be written
\[
\int B B\,\dd b=\frac{\dd\,\ell B}{\dd b},
\]
one may suppose that an expansion of the expression on the right will also lead to the goal. Since $B=\Gamma(b)\Gamma(1-b)$, it is enough to find an integral for $\frac{\dd\,\ell\Gamma(\mu)}{\dd\mu}$, which is known to be obtainable in the following way. From the definition of $\Gamma(\mu)$ it follows that
\[
\frac{\dd\,\ell\Gamma(\mu)}{\dd\mu}
=\frac1{\Gamma(\mu)}
  \int_0^\infty x^{\mu-1}e^{-x}\ell x\,\dd x.
\]
If for $\ell x$ one substitutes the integral
\[
\ell x=\int_0^\infty\frac{e^{-z}-e^{-zx}}{z}\dd z,
\]
which is obtained from the integral $\int_0^1y^{x-1}\dd y=\frac1x$ by integration with respect to $x$ between the limits $1$ and $x$, and by substituting $e^{-z}$ for $y$, then by reversing the order of integration and with the aid of equation (3) one very easily finds
\[
\frac{\dd\,\ell\Gamma(\mu)}{\dd\mu}
=\int_0^\infty\frac{\dd z}{z}
\left(e^{-z}-\frac1{(z+1)^\mu}\right).
\]
Putting $b$ and $(1-b)$ for $\mu$ gives
\[
\frac{\dd\,\ell B}{\dd b}
=\int_0^\infty\frac{\dd z}{z}
\left(\frac1{(z+1)^{1-b}}-\frac1{(z+1)^b}\right),
\]
and if one puts here $z=x-1$ or $z=\frac1x-1$,
\[
\frac{\dd\,\ell B}{\dd b}
=\int_1^\infty\frac{x^b-x^{1-b}}{x-1}\frac{\dd x}{x}
=\int_0^1\frac{x^b-x^{1-b}}{x-1}\frac{\dd x}{x}.
\]
Thus also
\[
2\frac{\dd\,\ell B}{\dd b}
=\int_0^\infty\frac{x^b-x^{1-b}}{x-1}\frac{\dd x}{x}.
\]
Comparison with equation (27) gives at once
\begin{equation}
2\frac{\dd\,\ell B}{\dd b}
=\int_{1-b}^{b} B B\,\dd b
=2\int_{\frac12}^{b} B B\,\dd b,
\tag{29}
\end{equation}
since, by Article 3, $B$ has the same values for $b=\frac12+b'$ and for $b=\frac12-b'$. Differentiating this equation with respect to $b$ gives again the differential equation treated in the preceding article.

\art{11}

Once the proofs just given have put beyond doubt the correctness of the equation
\[
B(r,1-r)=\int_0^\infty\frac{x^{r-1}\dd x}{x+1}
=\frac{\pi}{\sin r\pi},
\]
where $r$ denotes a positive proper fraction, it follows immediately from Article 2 that the Eulerian integrals of the first kind whose two arguments have a positive integer as their sum can actually be represented. One finds
\begin{equation}
B(m+r,n-r)=
\frac{(m+r-1)\cdots(1+r)r\,(n-1-r)\cdots(2-r)(1-r)}
     {(m+n-1)(m+n-2)\cdots2\cdot1}
\frac{\pi}{\sin r\pi}.
\tag{30}
\end{equation}
From this one can also draw an important converse conclusion about the Eulerian integrals of the first kind. Namely, if one puts $m=0$, $n=1$, $r=\frac12$, one finds
\begin{equation}
\Gamma\!\left(\frac12\right)
=\int_0^\infty e^{-x}\frac{\dd x}{\sqrt{x}}
=2\int_0^\infty e^{-xx}\dd x=\sqrt{\pi},
\tag{31}
\end{equation}
and consequently, by equation (4),
\begin{equation}
\Gamma\!\left(n+\frac12\right)
=(2n-1)(2n-3)\cdots5\cdot3\cdot1\,\frac{\sqrt{\pi}}{2^n}.
\tag{32}
\end{equation}

\art{12}

After the cases have been assembled in Articles 2 and 11 in which the Eulerian integrals of both kinds can be represented without the aid of new functions, I shall return once more, in conclusion, to the integral $B$, in order to develop a few further relations between it and other integrals. Among the different forms in which it appears, the following two are quite interesting:
\[
\int_{-\infty}^{+\infty}
\frac{e^{(2b-1)z}+e^{-(2b-1)z}}{e^z+e^{-z}}\dd z
\quad\text{and}\quad
2\int_0^{\frac{\pi}{2}}(\tg\varphi)^{2b-1}\dd\varphi
=2\int_0^{\frac{\pi}{2}}(\tg\varphi)^{1-2b}\dd\varphi .
\]
More important, however, are its generalizations by the introduction of new constants. Among these belongs the integral
\begin{equation}
\int_0^\infty\frac{x^{b-1}\dd x}{x+c}
=\frac{\pi c^{b-1}}{\sin b\pi},
\tag{33}
\end{equation}
where $c$ must be positive; yet it can subsequently be proved that $c$ may also be imaginary. Namely, if one multiplies numerator and denominator of the function $\frac{x^{b-1}}{x+i}$ by $(x-i)$, then the corresponding integral splits into two others, which can be reduced to the integral $B$ by the substitution $xx=y$. Thus one finds the validity of equation (33) for imaginary $c$. By differentiation and integration with respect to $c$, one can then again derive a series of other integrals.

If $c_1,c_2,\ldots,c_n$ are mutually distinct positive or imaginary constants, further if $0<b<n$, and if
\[
f(x)=(x+c_1)(x+c_2)\cdots(x+c_n),\qquad
f'(x)=\frac{\dd f(x)}{\dd x},
\]
then one also finds
\begin{equation}
\int_0^\infty\frac{x^{b-1}\dd x}{f(x)}
=\frac{\pi}{\sin b\pi}
 \sum_{k=1}^{n}\frac{c_k^{\,b-1}}{f'(-c_k)},
\tag{34}
\end{equation}
by decomposing $\frac{x^\beta}{f(x)}$ into partial fractions, where $\beta$ denotes the greatest integer contained in $b$.

A further advantage of the proof given in Article 10 is that it directly teaches how to determine a related class of integrals. From equations (27) and (29) it follows immediately that
\[
\int_0^\infty\frac{x^b-x^{\frac12}}{1-x}\frac{\dd x}{x}
=\pi\cot b\pi,
\]
and consequently also
\begin{equation}
\int_0^\infty\frac{x^a-x^b}{1-x}\frac{\dd x}{x}
=\pi(\cot a\pi-\cot b\pi).
\tag{35}
\end{equation}
From this one also obtains the integral denoted by $\varphi(b)$ in Article 3:
\[
\begin{aligned}
\int_0^\infty\frac{1-x^\mu}{1-x^\nu}x^{b-1}\dd x
&=\frac1\nu\int_0^\infty
  \frac{x^{\frac b\nu}-x^{\frac{\mu+b}{\nu}}}{1-x}\frac{\dd x}{x}\\
&=\frac{\pi}{\nu}
  \frac{\sin\frac{\mu\pi}{\nu}}
       {\sin\frac{b\pi}{\nu}\sin\frac{(\mu+b)\pi}{\nu}}\\
&=\frac{2\pi}{\nu}
  \frac{\sin\frac{\mu\pi}{\nu}}
       {\cos\frac{\mu\pi}{\nu}-\cos\frac{(\mu+2b)\pi}{\nu}},
\end{aligned}
\]
and its minimum:
\[
\int_0^\infty
\frac{x^{\frac{\mu}{2}}-x^{-\frac{\mu}{2}}}
     {x^{\frac{\nu}{2}}-x^{-\frac{\nu}{2}}}
\frac{\dd x}{x}
=\frac{2\pi}{\nu}\,\tg\frac{\mu\pi}{2\nu}.
\]

Most of the integrals belonging here are found in the “Sammlung von Integraltafeln” issued by Minding at the commission of the Prussian ministry (Berlin, 1849), at the beginning of the fourth division. Perhaps this is the place to indicate and correct several errors found there. By decomposing $\frac1{(x-1)(x+c)}$ into partial fractions and applying the formulas of this article, one easily finds
\[
\int_0^\infty\frac{x^a-1}{(x-1)(x+c)}\dd x
=\frac{\pi}{c+1}
\left(\frac{c^a-\cos a\pi}{\sin a\pi}
-\frac{\ell c}{\pi}\right).
\]
This equation holds for positive and negative proper fractional $a$, as one easily verifies by writing $\frac1x$ and $\frac1c$ instead of $x$ and $c$. If one differentiates with respect to $a$ and then puts $a=0$, one obtains a series of integrals that are given incorrectly in those tables. To make the tedious differentiations as easy as possible, the following device may be used. Put
\[
\frac{\ell c}{\pi}+\frac{c+1}{\pi}
\int_0^\infty\frac{x^a-1}{(x-1)(x+c)}\dd x
=\frac{c^a-\cos a\pi}{\sin a\pi}=J,
\qquad
\frac{\dd^nJ}{\dd a^n}={}^{n}J.
\]
Then for $a=0$ one obtains
\[
\int_0^\infty\frac{(\ell x)^n\dd x}{(x-1)(x+c)}
=\frac{\pi}{c+1}\,{}^{n}J_0.
\]
But one easily finds
\[
\begin{aligned}
&J\pi^n\sin\left(a+\frac n2\right)\pi
+n\,{}^{1}J\,\pi^{n-1}\sin\left(a+\frac{n-1}{2}\right)\pi+\cdots\\
&\quad{}+n\,{}^{n-1}J\,\pi\sin\left(a+\frac12\right)\pi
+{}^{n}J\sin a\pi\\
&=\frac{\dd^n(J\sin a\pi)}{\dd a^n}
=c^a(\ell c)^n-\pi^n\cos\left(a+\frac n2\right)\pi.
\end{aligned}
\]
For $a=0$ this gives recursion formulas for the $J_0$ with even indices and those with odd indices. Thus one finds
\[
\begin{aligned}
\int_0^\infty\frac{\ell x\,\dd x}{(x-1)(x+c)}
&=\frac{(\ell c)^2+\pi^2}{2(c+1)},\\
\int_0^\infty\frac{(\ell x)^2\dd x}{(x-1)(x+c)}
&=\frac{\ell c\bigl((\ell c)^2+\pi^2\bigr)}{3(c+1)},\\
\int_0^\infty\frac{(\ell x)^3\dd x}{(x-1)(x+c)}
&=\frac{\bigl((\ell c)^2+\pi^2\bigr)^2}{4(c+1)},\\
\int_0^\infty\frac{(\ell x)^4\dd x}{(x-1)(x+c)}
&=\frac{\ell c\bigl((\ell c)^2+\pi^2\bigr)
        \bigl(3(\ell c)^2+7\pi^2\bigr)}{5(c+1)\cdot3},\\
\int_0^\infty\frac{(\ell x)^5\dd x}{(x-1)(x+c)}
&=\frac{\bigl((\ell c)^2+\pi^2\bigr)^2
        \bigl((\ell c)^2+3\pi^2\bigr)}{6(c+1)}.
\end{aligned}
\]
In those tables, however, in place of the divisors $2,3,4,5,6$, the divisors $1\cdot2$, $1\cdot2\cdot3$, $1\cdot2\cdot3\cdot4$, $1\cdot2\cdot3\cdot4\cdot5$, $1\cdot2\cdot3\cdot4\cdot5\cdot6$ are given.

\end{document}
